% simple-minimizer-existence.tex — \input'd from proofs.tex
% Conventions: .claude/rules/thesis-tex.md
%
% Sources: Jörn's dictation + Jörn's 120-minute talk (January 2025) expanding
%   HK2017 proof. Notation per appendix-notation.tex.
%
% Structure: Simple orbit definitions, existence theorems, proof via primal-dual
%   and simplification algorithm

Next, we use the established dual problem to prove the existence of "simple" minimizers, which unlike arbitrary generalized Reeb orbits, have a finite dimensional combinatorial structure, and are thus computationally representable, and tractable as search spaces for optimization problems.
The existence of simple minimizers has been considered surprising by the authors, until encountering the proof in \cite{HK2017} and fleshing it out into the proof below in this thesis.
There is a-priori no reason to expect that for the rigid Reeb orbits, which can be isolated critical points of the primal problem, there would be a construction that changes their combinatorial structure and yet preserves minimality.

\begin{definition}[Simple Reeb orbit]\label{def:simple-reeb-orbit}
% Jörn: text approved (05aa8b3)
% QC: "simple" is HK2017's terminology ("simple closed characteristic").
%   Our definition uses the Reeb parametrization (gamma_dot = R_i exactly,
%   not gamma_dot = C J n_i with varying C > 0 as in HK2017). Equivalent
%   up to reparametrization; our gamma_dot is Reeb-normalized (lambda_0 = 1).
%   Key properties: (1) piecewise constant velocity, (2) pure facet Reeb
%   vectors (not convex combinations), (3) each facet used at most once.
%   Contrast def:reeb-orbit-polytope which allows convex combinations.
% Downstream: simple orbits are the computational objects. The algorithm
%   searches over all simple orbits by enumerating (sigma, beta).
%   Rust: stored in EhzResult::best_permutation (Vec<usize>) for σ.
%   Dwell times τ are implicit in best_beta; no SimpleOrbit struct exists.
A generalized Reeb orbit $\gamma$ on a polytope $K$
(Definition~\ref{def:reeb-orbit-polytope}) is \emph{simple} if:
\begin{enumerate}
\item $\dot{\gamma}$ is piecewise constant;
\item each constant value is a facet Reeb vector $R_i$
  (not a nontrivial convex combination); and
\item for each facet index $i$, the set
  $\{ t \in [0, T] : \dot{\gamma}(t) = R_i \}$
  is a single interval (possibly empty; in a non-simple orbit,
  a facet may contribute multiple or infinitely many disjoint intervals).
\end{enumerate}
\end{definition}

\begin{remark}[Finite representation of simple orbits]\label{rem:simple-orbit-data}
% Jörn: text approved (05aa8b3)
% QC: breakpoints are velocity changes, not facet transitions.
%   Example (Jörn): orbit on a Lagrangian 2-face F_1 cap F_2,
%   traveling first with R_1, breakpoint in int(F_1 cap F_2),
%   then R_2. Facet incidence {F_1, F_2} is constant across the breakpoint.
% Downstream: Rust stores σ in EhzResult::best_permutation: Vec<usize>
%   (0-based cyclic permutation of facet indices). Dwell times τ are not
%   stored separately; they are encoded in best_beta: Vec<f64> (β_k ∝ τ_k,
%   normalised via η^Tβ = 1). No standalone SimpleOrbit struct exists.
A simple Reeb orbit can be described by a triple $(\sigma, \tau, b)$:
\begin{itemize}
\item a permutation $\sigma$ of $\{1, \ldots, F\}$
  (the cyclic order in which facet Reeb vectors appear as velocities),
\item durations $\tau_1, \ldots, \tau_F \geq 0$
  (where $\tau_k$ is the time spent with velocity $R_{\sigma(k)}$;
  $\tau_k = 0$ means facet $\sigma(k)$ is not visited), and
\item a base point $b = \gamma(0) \in \partial K$.
\end{itemize}
Not every triple $(\sigma, \tau, b)$ yields a valid orbit;
the orbit starting at~$b$ with the prescribed velocities must
remain on~$\partial K$, and closure requires
$\sum_k \tau_k R_{\sigma(k)} = 0$.
The period is $T = \sum_k \tau_k$ and the action is $A(\gamma) = T$
(Remark~\ref{rem:contact-normalization}).
By Lemma~\ref{lem:shoelace}, the action depends only on
$(\sigma, \tau)$, not on the base point~$b$.

Two triples $(\sigma, \tau, b)$ and $(\sigma', \tau', b')$
represent the same orbit if and only if, after removing all entries
with $\tau_k = 0$, the resulting subsequences of (facet, duration)
pairs agree up to cyclic shift and $b'$ is the corresponding
point on the orbit.

Breakpoints of $\dot{\gamma}$ are velocity changes, not necessarily
facet transitions: for example, the orbit may travel along a
lower-dimensional face $F_i \cap F_j$ with velocity $R_i$, then
switch to $R_j$ at a breakpoint in the interior of $F_i \cap F_j$.
\end{remark}

% Jörn: math approved (f7d26ea) — from \begin{lemma}[Base point recovery] to \end{remark} (rem:beta-to-tau)

\begin{lemma}[Base point recovery]\label{lem:base-point-recovery}
% QC: the tangency condition ⟨n_i, R_i⟩ = 0 (Remark~\ref{rem:reeb-properties})
%   eliminates time-dependence from the on-facet constraint, reducing it to a
%   linear system in b. The proof is a direct consequence of this tangency.
% Downstream: algorithm for computing b from (σ, β) output of Algorithm~\ref{alg:ehz}.
%   Rust: recover_base_point() in algorithms/hk2017/recover.rs.
Let $K$ be a polytope and $(\sigma, \tau, b)$ a simple orbit
(Definition~\ref{def:simple-reeb-orbit}) with $m$ active facets
$S = \{\sigma(k) : \tau_k > 0\}$.
Define the accumulated displacements
\[
  v_0 = 0, \qquad
  v_k = \sum_{j=1}^{k-1} \tau_j\, R_{\sigma(j)}
  \quad \text{for } k = 1, \ldots, m.
\]
A point $b \in \mathbb{R}^4$ is a valid base point
(i.e., the orbit $\gamma$ with $\gamma(0) = b$ and
velocity $R_{\sigma(k)}$ on the $k$-th segment
is a generalized Reeb orbit on~$\partial K$) if and only if:
\begin{enumerate}
\item \emph{On-facet (equality):}
  for each active $k$,
  \begin{equation}\label{eq:base-point-equality}
    \langle n_{\sigma(k)},\, b \rangle
    = h_{\sigma(k)} - \langle n_{\sigma(k)},\, v_k \rangle;
  \end{equation}
\item \emph{Inside $K$ (inequality):}
  for each active segment $k$ and all $s \in [0, \tau_k]$,
  \begin{equation}\label{eq:base-point-inequality}
    \langle n_j,\, b + v_k + s\, R_{\sigma(k)} \rangle
    \leq h_j
    \quad \text{for all } j = 1, \ldots, F.
  \end{equation}
\end{enumerate}
The equality constraints~\eqref{eq:base-point-equality} form
a linear system $N_S\, b = r$ with $m$ equations and $4$~unknowns,
where $N_S$ is the $m \times 4$ matrix with rows $n_{\sigma(k)}^\top$
and $r_k = h_{\sigma(k)} - \langle n_{\sigma(k)}, v_k \rangle$.
The set of valid base points is a convex polytope of
dimension at most $4 - \operatorname{rank}(N_S)$.
\end{lemma}

\begin{proof}
\emph{Equality constraints.}
During the $k$-th segment, the orbit position is
\[
  \gamma(t) = b + v_k + (t - t_k)\, R_{\sigma(k)},
  \quad t \in [t_k,\, t_k + \tau_k],
\]
where $t_k = \sum_{j=1}^{k-1} \tau_j$.
For $\gamma$ to lie on facet $F_{\sigma(k)}$,
we require $\langle n_{\sigma(k)},\, \gamma(t) \rangle
= h_{\sigma(k)}$ throughout the segment.
Since $R_{\sigma(k)}$ is tangent to $F_{\sigma(k)}$
($\langle n_{\sigma(k)}, R_{\sigma(k)} \rangle = 0$
by Remark~\ref{rem:reeb-properties}),
the time-dependent term vanishes:
\begin{align*}
  \langle n_{\sigma(k)},\, \gamma(t) \rangle
  &= \langle n_{\sigma(k)},\, b + v_k \rangle
  + (t - t_k)\,
    \langle n_{\sigma(k)},\, R_{\sigma(k)} \rangle \\
  &= \langle n_{\sigma(k)},\, b + v_k \rangle.
\end{align*}
Hence the on-facet condition for the entire segment
reduces to the single
equation~\eqref{eq:base-point-equality}.

\emph{Inequality constraints.}
The orbit must remain inside~$K$:
$\langle n_j, \gamma(t) \rangle \leq h_j$
for all $j$ and all~$t$.
Since $\gamma(t)$ is affine in~$t$ on each segment,
each inequality $\langle n_j, b + v_k + s\, R_{\sigma(k)} \rangle
\leq h_j$ is linear in~$s$.
A linear function on $[0, \tau_k]$ attains its
maximum at an endpoint, so it suffices to
check~\eqref{eq:base-point-inequality}
at $s = 0$ and $s = \tau_k$ only.
These are precisely the breakpoints
$b + v_k$ and $b + v_{k+1}$.

\emph{Solution structure.}
The equality constraints are a linear system
$N_S\, b = r$ in~$b \in \mathbb{R}^4$.
Its solution set is an affine subspace of
dimension $4 - \operatorname{rank}(N_S)$.
The inequality constraints at the $2m$ breakpoints
(with $s \in \{0, \tau_k\}$) are linear in~$b$,
so their feasible set is a polyhedron.
The intersection of the affine subspace with
the polyhedron is a convex polytope (bounded,
since $K$ is bounded).
\end{proof}

\begin{remark}[Non-uniqueness and computation]\label{rem:base-point-computation}
% Downstream: implementation uses SVD least-norm solution.
%   When rank(N_S) < 4, picks the minimum-norm b in the affine subspace.
%   Inequality constraints are checked post-hoc.
The solution space has dimension $4 - \operatorname{rank}(N_S)$.
If $\operatorname{rank}(N_S) = 4$, then $b$ is unique;
otherwise the inequality constraints select a
convex subset of the affine solution space.
Generically, orbits visiting $\geq 4$ facets have
$\operatorname{rank}(N_S) = 4$; rank deficiency occurs
when active normals are linearly dependent
(e.g., product polytopes with parallel normal pairs).
The capacity does not depend on~$b$
(Lemma~\ref{lem:shoelace}).
In practice, we compute the least-norm solution
of $N_S\, b = r$ via SVD and verify the
inequality constraints.\end{remark}

\begin{remark}[Conversion from algorithm output]\label{rem:beta-to-tau}
% QC: connects the KKT output (σ, β) to the dwell times τ used above.
%   The formula τ_k = T h_{σ(k)} β_{σ(k)} comes from the variable
%   substitution in the proof of Theorem~\ref{thm:main}.
Algorithm~\ref{alg:ehz} outputs a permutation~$\sigma$ and
a vector~$\beta$ with $\beta_i \geq 0$ and
$\sum_i \beta_i h_i = 1$.
The dwell times are
\begin{equation}\label{eq:beta-to-tau}
  \tau_k = T\, h_{\sigma(k)}\, \beta_{\sigma(k)},
\end{equation}
where $T = c_{\mathrm{EHZ}}(K)$ is the capacity
(the period of the orbit).
Substituting into Lemma~\ref{lem:base-point-recovery}
recovers a base point from the algorithm output.
\end{remark}

\begin{theorem}[Existence of closed characteristics --- smooth case]\label{thm:orbit-existence-smooth}
% Jörn: text approved (05aa8b3)
% QC: classical result, Rabinowitz 1978.
% Downstream: foundation for the EHZ capacity on smooth convex bodies.
Let $K$ be a convex body in $\mathbb{R}^4$ with smooth boundary~\cite{Rabinowitz1978}.
There exists a closed characteristic on $\partial K$.
The infimum of the action over all closed characteristics
is attained and positive:
\[
  c_{\mathrm{EHZ}}(K) = \min_{\gamma} A(\gamma) > 0.
\]
\end{theorem}

\begin{theorem}[Existence of closed characteristics --- polytope case]\label{thm:orbit-existence-polytope}
% Jörn: text approved (05aa8b3)
% QC: Artstein-Avidan & Ostrover 2014 extend to general convex bodies
%   (including polytopes) via generalized Reeb orbits.
% Downstream: justifies c_EHZ being well-defined for polytopes.
Let $K$ be a polytope in $\mathbb{R}^4$~\cite{AAO2014}.
There exists a generalized Reeb orbit on $\partial K$
(Definition~\ref{def:reeb-orbit-polytope}).
The minimum action over all generalized Reeb orbits is attained
and positive:
\[
  c_{\mathrm{EHZ}}(K) = \min_{\gamma} A(\gamma) > 0.
\]
\end{theorem}

\begin{theorem}[Simple minimizer {\cite{HK2017}}]\label{thm:simple-minimizer}
% Jörn: text approved (05aa8b3)
% QC: main theorem of Haim-Kislev 2017. Deep result; we cite, not prove.
%   The proof constructs a minimizing sequence and shows it can be
%   "simplified" (breaking convex combinations into pure Reeb vectors,
%   merging repeated facets into contiguous intervals).
%   Jörn gave a 120-minute talk on this proof.
% Downstream: reduces the infinite-dimensional optimization (min over all
%   generalized Reeb orbits) to a finite search (over simple orbits).
%   Combined with the finite representation (rem:simple-orbit-data),
%   this makes the problem computationally tractable.
Let $K$ be a polytope.
Among all generalized Reeb orbits on $\partial K$, there exists one
with minimum action that is simple
(Definition~\ref{def:simple-reeb-orbit}).
Consequently,
\[
  c_{\mathrm{EHZ}}(K) = \min\bigl\{ A(\gamma) :
  \gamma \text{ is a simple Reeb orbit on } \partial K \bigr\}.
\]
\end{theorem}



% Jörn: text approved (4d6f7a6) — 1:1 restatement of thm:simple-minimizer
\begin{theorem*}[Theorem~\ref{thm:simple-minimizer} restated]
Let $K$ be a polytope.
Among all generalized Reeb orbits on $\partial K$, there exists one
with minimum action that is simple
(Definition~\ref{def:simple-reeb-orbit}).
Consequently,
\[
  c_{\mathrm{EHZ}}(K) = \min\bigl\{ A(\gamma) :
  \gamma \text{ is a simple Reeb orbit on } \partial K \bigr\}.
\]
\end{theorem*}

The proof constructs a simple dual minimizer from an arbitrary one
via five operations, each captured by a standalone lemma.
We state all five lemmas, then prove the theorem by chaining them,
and finally prove each lemma.

% Jörn: text approved (6f31407) — lemma statements 5.10–5.14 and theorem proof

% QC: translates HK2017 Lemma 4.2.
%   Core: W^{1,2} approximation preserving closure and velocity constraint.
%   Corollary: A and I_K estimates by continuity.
\begin{lemma}[Piecewise affine approximation]\label{lem:pw-affine-approx}
Let $z \in W^{1,2}([0,T], \mathbb{R}^4)$ be closed
with $\dot{z}(t) \in \mathrm{conv}\{R_1, \ldots, R_F\}$ a.e.
For any $\varepsilon > 0$, there exists a closed
piecewise affine curve~$z'$ with
$\dot{z}'(t) \in \mathrm{conv}\{R_1, \ldots, R_F\}$ on each piece
and $\|z' - z\|_{W^{1,2}} < \varepsilon$.

In particular,
$|A(z') - A(z)| < C\varepsilon$ and
$|I_K(z') - I_K(z)| < C\varepsilon$
for a constant~$C$ depending only on~$z$.
\end{lemma}

% QC: translates HK2017 Lemma 4.4.
%   Core: split convex-combination velocities into pure Reeb, A non-decreasing.
%   I_K = T is an immediate consequence, not the core difficulty.
\begin{lemma}[Splitting mixed velocities]\label{lem:splitting}
Let $z \colon [0,T] \to \mathbb{R}^4$ be a closed
piecewise affine curve with velocities in
$\mathrm{conv}\{R_1, \ldots, R_F\}$.
There exists a closed piecewise affine curve~$z'$
with the same period~$T$, whose velocity on each segment is a single
Reeb vector~$R_i$, such that
$A(z') \geq A(z)$.
In particular, $I_K(z') = T$.
\end{lemma}

% QC: translates HK2017 Lemma 4.5.
\begin{lemma}[Merging repeated velocities]\label{lem:merging}
Let $z \colon [0,T] \to \mathbb{R}^4$ be a closed piecewise affine
curve whose velocity on each segment is a single Reeb vector~$R_i$.
If some~$R_k$ appears in two non-contiguous segments, there
exists a rearrangement~$z'$ that merges them into one contiguous
segment with $A(z') \geq A(z)$ and $I_K(z') = I_K(z)$.
Iterating until each Reeb vector appears at most once preserves
these bounds.
\end{lemma}

% QC: immediate from shoelace (A quadratic in durations) and h_K^2 normalization (I_K linear).
\begin{lemma}[Rescaling]\label{lem:rescaling}
Let $z \colon [0,T] \to \mathbb{R}^4$ be a closed piecewise affine
curve with pure Reeb velocities.
For any $\beta > 0$, multiplying all durations by~$\beta$
gives a curve~$z'$ with
$A(z') = \beta^2 A(z)$ and $I_K(z') = \beta\, I_K(z)$.
\end{lemma}

% QC: compactness of S_F × Δ^{F-1} + continuity of A in (σ, s, T).
%   Handles varying periods T_n → T.
% Jörn: text approved (6f31407)
\begin{lemma}[Compactness of simple curves]\label{lem:compactness}
Let $(z_n)$ be a sequence of closed simple
piecewise-Reeb curves with periods $T_n \to T > 0$.
Then there are translation vectors $b_n$ such that a subsequence of
$(z_n + b_n)$ converges to a closed simple piecewise-Reeb curve~$z_*$
with period~$T$, $A(z_*) = \lim A(z_n)$, and $I_K(z_*) = T$.
\end{lemma}

%%% --- Proof of the theorem ---

\begin{proof}[Proof of Theorem~\ref{thm:simple-minimizer}]
A minimum-action generalized Reeb orbit on~$\partial K$ exists
(Theorem~\ref{thm:orbit-existence-polytope}).
By the primal-dual equivalence
(Theorem~\ref{thm:primal-dual-equivalence}),
any primal minimizer~$\gamma$ is also a dual minimizer,
say $z_0 := \gamma$, with
$I_K(z_0) = T = A(\gamma)$.
Since $\gamma$ is a generalized Reeb orbit on a polytope,
$\dot{z}_0(t) = \dot{\gamma}(t)
  \in \mathrm{conv}\{R_1, \ldots, R_F\}$ a.e.

For each $\varepsilon > 0$, apply the first four lemmas in sequence:
\begin{enumerate}
\item \emph{Approximate}
  (Lemma~\ref{lem:pw-affine-approx}):
  obtain a closed piecewise affine $z_1$
  with velocities in $\mathrm{conv}\{R_i\}$,
  $|A(z_1) - T| < \varepsilon$, and
  $|I_K(z_1) - T| < \varepsilon$.

\item \emph{Split}
  (Lemma~\ref{lem:splitting}):
  obtain a closed $z_2$ with pure Reeb velocities,
  $A(z_2) \geq A(z_1)$, and $I_K(z_2) = T$.

\item \emph{Merge}
  (Lemma~\ref{lem:merging}):
  obtain a closed $z_3$ with each Reeb vector
  used at most once,
  $A(z_3) \geq A(z_2)$, and $I_K(z_3) = T$.

\item \emph{Rescale}
  (Lemma~\ref{lem:rescaling}):
  set $\beta = T / A(z_3)$ and obtain $z_4$ with
  period $T_4 = \beta\, T = T^2 / A(z_3)$,
  $A(z_4) = \beta^2 A(z_3) = T_4$, and
  $I_K(z_4) = \beta\, T = T_4$.
\end{enumerate}

\emph{Optimality bound.}
The curve $z_4$ is feasible for the dual problem
with period~$T_4$
(closed, $A(z_4) = T_4$).
Since $z_0$ was a dual minimizer, the dual-feasible
curve~$z_4$ satisfies
$I_K(z_4) = T^2 / A(z_3) \geq I_K(z_0) = T$,
i.e., $A(z_3) \leq T$.
Combined with $A(z_3) \geq A(z_1) \geq T - \varepsilon$,
we get $A(z_3) \to T$ as $\varepsilon \to 0$,
hence $\beta \to 1$, $I_K(z_4) \to T$, and $T_4 \to T$.

\begin{enumerate}\setcounter{enumi}{4}
\item \emph{Compactness}
  (Lemma~\ref{lem:compactness}):
  the sequence $(z_4)$ (indexed by $\varepsilon \to 0$)
  has periods $T_4 \to T$, so a subsequence converges up to translation
  to a closed simple piecewise-Reeb
  curve~$z_*$ with
  $A(z_*) = T$ and $I_K(z_*) = T$.
  Thus $z_*$ is a dual minimizer.
\end{enumerate}

By Theorem~\ref{thm:primal-dual-equivalence},
$z_*$ corresponds to a primal minimizer
$\gamma_* = \nu z_* + b$ for some constant
$\nu > 0$, $b \in \mathbb{R}^4$.
Since $z_*$ is simple, so is~$\gamma_*$.
\end{proof}

%%% --- Proofs of the lemmas ---
% Jörn: text approved (6f31407) — proofs of lemmas 5.10–5.14

\begin{proof}[Proof of Lemma~\ref{lem:pw-affine-approx}]
% QC: standard density argument + integral mean value in a convex set.
Since $\dot{z}(t) \in \mathrm{conv}\{R_1, \ldots, R_F\}$ a.e.,
the velocity~$\dot{z}$ is bounded a.e.
Hence we can find a partition $0 = t_0 < t_1 < \cdots < t_m = T$
fine enough that the piecewise affine interpolant~$z'$
defined by $z'(t_k) = z(t_k)$ (affine on each
$(t_k, t_{k+1})$) satisfies
$\|z' - z\|_{L^\infty} < \varepsilon$ and
$\|\dot{z}' - \dot{z}\|_{L^\infty} < \varepsilon$.

On each subinterval $(t_k, t_{k+1})$, the velocity
of~$z'$ is the integral mean of~$\dot{z}$:
\[
  \dot{z}'(t)
  = \frac{z(t_{k+1}) - z(t_k)}{t_{k+1} - t_k}
  = \frac{1}{t_{k+1} - t_k}
    \int_{t_k}^{t_{k+1}} \dot{z}(s)\, ds.
\]
Since $\dot{z}(s) \in \mathrm{conv}\{R_1, \ldots, R_F\}$
a.e.\ and the convex hull is closed, the integral mean
lies in $\mathrm{conv}\{R_1, \ldots, R_F\}$.

Closure: $z'(0) = z(0) = z(T) = z'(T)$.

The ``in particular'' clause follows from
continuity of $A$ and $I_K$ in the $W^{1,2}$ norm:
$A$ is bilinear in $(z, \dot{z})$, and
$I_K$ is continuous in~$\dot{z}$ since $h_K^2$ is
locally Lipschitz.
\end{proof}

\begin{proof}[Proof of Lemma~\ref{lem:splitting}]
% QC: core idea: split each mixed-velocity segment into pure segments,
%   choosing the ordering to preserve or increase action.
%   Then show I_K = T by direct computation from Reeb normalization.
Process each segment independently.
If segment~$I$ has velocity
$v = \sum_{i=1}^k a_i R_i$ with $a_i > 0$,
$\sum a_i = 1$,
replace~$I$ by $k$ consecutive subsegments of
lengths $\Delta t_i = a_i |I|$, each with pure
velocity~$R_i$.

\emph{Closure preserved:}
the total displacement is
$\sum_i \Delta t_i R_i
  = |I| \sum_i a_i R_i = |I|\, v$,
the same as the original segment.

\emph{Period preserved:}
$\sum_i \Delta t_i = |I| \sum_i a_i = |I|$.

\emph{Action non-decrease:}
by Lemma~\ref{lem:shoelace},
splitting segment~$I$ introduces
new internal cross-terms
$\sum_{r < s} \Delta t_r\, \Delta t_s\,
  \omega_0(R_r, R_s)$.
Reversing the ordering of the $k$ subsegments
flips the sign of this sum
(since $\omega_0(R_r, R_s) = -\omega_0(R_s, R_r)$).
Hence one of the two orderings yields
$\Delta A \geq 0$.

\emph{Dual functional:}
for a pure Reeb velocity
$R_i = \frac{2}{h_i} J_0 n_i$,
\[
  h_K^2(-J_0 R_i)
  = h_K^2\Bigl(\frac{2}{h_i} n_i\Bigr)
  = \frac{4}{h_i^2}\, h_K^2(n_i)
  = \frac{4}{h_i^2} \cdot h_i^2
  = 4,
\]
using $h_K(n_i) = h_i$
(Definition~\ref{def:facets}).
Hence
$I_K(z') = \frac{1}{4} \int_0^T 4\, dt = T$.
\end{proof}

\begin{proof}[Proof of Lemma~\ref{lem:merging}]
% QC: the forward/backward merge argument uses the action identity.
%   Key: the two action differences have opposite signs,
%   so one is always non-negative.
Let $w_1, \ldots, w_m$ be the velocity sequence with
$w_r = w_s = v$ for some $r < s$.
Write $M = \{r{+}1, \ldots, s{-}1\}$ for the middle segments.

\emph{Forward merge} (merge to position~$r$):
remove segment~$I_s$ and extend the duration of~$I_r$
by~$|I_s|$,
placing the middle segments after the merged segment.
By Lemma~\ref{lem:shoelace}:
\[
  2A(z) = \sum_{1 \leq j < i \leq m}
    |I_j|\,|I_i|\,\omega_0(w_j, w_i).
\]
The rearrangement only affects cross-terms
between the middle segments~$M$ and the two
segments with velocity~$v$ at positions $r$, $s$.
The action difference between the merged and
original curves is:
\[
  \Delta A_{\mathrm{fwd}}
  = |I_s| \sum_{i \in M} |I_i|\, \omega_0(v, w_i).
\]

\emph{Backward merge} (merge to position~$s$):
extend $I_s$ and place the middle segments before it.
The same cross-term calculation gives:
\[
  \Delta A_{\mathrm{bwd}}
  = -|I_r| \sum_{i \in M} |I_i|\, \omega_0(v, w_i).
\]

Since $|I_r|, |I_s| > 0$, the two action changes have
opposite signs (or both vanish).
Hence
$\max(\Delta A_{\mathrm{fwd}},
  \Delta A_{\mathrm{bwd}}) \geq 0$.

Closure is preserved: the total displacement
$\sum |I_k| w_k$ is unchanged by rearrangement.
The dual functional $I_K$ is unchanged:
$h_K^2(-J_0 \dot{z})$ depends only on the velocity
at each time, not on its position in the sequence,
so reordering segments preserves~$\int h_K^2$.
\end{proof}

\begin{proof}[Proof of Lemma~\ref{lem:rescaling}]
By Lemma~\ref{lem:shoelace}, every term in $2A(z)$
contains the product $|I_i|\,|I_j|$, so
$A(z') = \beta^2 A(z)$.
Since $h_K^2(-J_0 \dot{z})$ depends only on the velocity
(unchanged by rescaling), $I_K$ scales linearly:
$I_K(z') = \beta\, I_K(z)$.
\end{proof}

\begin{proof}[Proof of Lemma~\ref{lem:compactness}]
% Jörn: text approved (6f31407)
% QC note: the encoding (sigma, s, T) determines a piecewise-Reeb curve
%   up to translation. Translation is a free parameter for dual curves.
%   Compactness requires choosing appropriate translations b_n to ensure
%   the sequence stays bounded.
%   (Contrast Remark~\ref{rem:simple-orbit-data}, where primal orbits gamma
%   need a base point b = gamma(0) in addition to (sigma, tau).)
Each $z_n$ is encoded by a tuple
$(\sigma_n, s_n, T_n, z_n(0))$
in $S_F \times \Delta^{F-1} \times (0,\infty) \times \mathbb{R}^4$,
where $\sigma_n$ is the ordering,
$s_{n,i} = t_{n,i} / T_n$ are normalized durations,
and $z_n(0)$ is the start point needed to trace the piecewise affine path.
Since $T_n \to T > 0$, eventually
$T_n \in [T/2,\, 2T]$.
We only want convergence up to translations, so we drop the $z_n(0)$.
The space
$S_F \times \Delta^{F-1} \times [T/2,\, 2T]$
is compact, so a subsequence has
$(\sigma_n, s_n, T_n) \to (\sigma_*, s_*, T)$.
The curve~$z_*$ with parameters
$(\sigma_*, s_*, T)$ is the limit (up to choice of starting point).
By Lemma~\ref{lem:shoelace},
$A(\sigma, s, T) = T^2 \sum_{j<i} s_i\, s_j\,
\omega_0(w_j, w_i)$ is continuous in $(\sigma, s, T)$,
so $A(z_*) = \lim A(z_n)$.
Since all velocities are pure Reeb vectors,
$I_K(z_*) = T$.
\end{proof}

